\documentclass{article}

\usepackage[margin=1in]{geometry}

\usepackage{parskip}

\begin{document}

% Page 1

This journal is ©The Royal Society of Chemistry 2021
Chem. Commun., 2021, 57, 595--598 | 595
Cite this: Chem. Commun., 2021,
57, 595
Low-spin 1,10-diphosphametallocenates of
chromium and iron†
Samuel M. Greer,
abc O¨ kten U¨ngor,
b Ross J. Beattie,
a
Jaqueline L. Kiplinger,
*a Brian L. Scott,
a Benjamin W. Stein
*a and
Conrad A. P. Goodwin
*ad
We report two anionic diphosphametallocenates, [K(2.2.2-crypt)]-
[M(PC4Me4)2] (M = Cr, 2-Cr; Fe, 2-Fe). Both are low-spin (S = 1
2) by EPR
spectroscopy and SQUID magnetometry. This contrasts the high-spin
(S = 3
2) ferrocenate, [K(2.2.2-crypt)][Fe(C5H2-1,2,4-tBu)2] (4-Fe). Quantum
chemical calculations suggest this is due to significant differences in
ligand field splitting of the d-orbitals which also explain structural
features in the 2-M complexes.
The chemistry and electronic structure of transition metal (TM)
metallocenes1,2 has fascinated chemists since the discovery of
ferrocene, [Fe(Cp)2] (FcH, Cp = cyclopentadienyl, \{C5H5\}).3
Oxidation of FcH leads to the ferrocenium cation, [Fe(Cp)2]+
(FcH+),4 which finds use in organic synthesis,5 while the FcH+/0
couple
serves
as
an
electrochemical
standard.6
‘‘Hetero-
metallocenes’’ are related complexes that feature heterocycles in
lieu of Cp ligands (e.g. Z5-pyrrolides, \{NC4R4\}; Z5-stannolides
\{SbC4R4\}; poly-phospholides \{PnC5nR5n\} etc.).7 The exchange
of C atoms for heteroatoms can lead to drastic changes in
reactivity.1,2 For example, ferrocene undergoes deprotonation by
alkyl-lithium reagents,8,9 but 1,10-diphosphaferrocenes undergo
nucleophilic attack at the P atom.10,11
The oxidation of TM metallocenes to metallocenium cations
is a defining feature of their chemistry, with examples across the
first row transition metals (V–Ni).1,12 The corresponding reduc-
tive chemistry is underdeveloped outside of electrochemical
studies.13 Exceptions include the synthesis of [Mn(Cp*)2]
salts,14 the reduction of [Co\{Z5-C9H5-1,3-(SiMe3)2\}2] to a formally
anionic ‘‘cobaltate’’,15 and an anionic iron bis-stannolide.7c
Some of us have recently described the chemical reduction of
[M(Cpttt)2] (3-M; M = Mn, Fe, Co; Cpttt = \{C5H2-1,2,4-tBu\}) to
[K(2.2.2-crypt)][M(Cpttt)2] (4-M; M = Mn, Fe, Co) giving an 18,
19 and 20 valence electron (VE) series,16 similar to the 17, 18, 19
VE series of [Fe(CpR)(arene)]2+/+/0 complexes.17
The synthesis of 3d metallocene anions could open up the
possibility of using them as
nucleophiles like ‘‘A[Re(Cp)2]’’
(A = Li,18 K19). While the [Mn(Cp*)2] anion is temperature
stable,14c,a analogous [FcH] and [Fe(Cp*)2] anions are not easily
accessed chemically;13b and 4-M complexes are extremely tempera-
ture sensitive.16 We sought ligands that would impart greater
stability to the anionic species. Reports of [M(PC4RnH4n)2] for a
range of metals include the Fe/Ru/Os triad,7d,20 and many exhibit
facile redox chemistry.21 For example, [M(TMP)2] (M = Cr, 1-Cr;22 Fe,
1-Fe;23 TMP = \{PC4Me4\}) have M+/0 couples (vs. FcH+/0: M = Cr,
0.77 V;24 Fe, 0.06 V), and M0/ (vs. FcH+/0: M = Cr, 2.4 V; Fe,
3.02 V) at less cathodic potentials than 3-M.16 Here we report the
synthesis of 1,10-diphosphametallocenates, [K(2.2.2-crypt)][M(TMP)2]
(M = Cr, 2-Cr; Fe, 2-Fe), by low-temperature reduction of 1-M by KC8
in the presence of 2.2.2-cryptand (Scheme 1).
Precursor 1-M (M = Cr, Fe) were synthesized from MCl2 and
base-free KTMP by modification of literature procedures.21b,23b
During the course of this work we have also structurally
characterized [Zr(Cp)2(C4Me4)] (5) and Ph-TMP (6) for the first
time, which are precursors in the synthesis of KTMP (see ESI†).
Complexes
1-M
were
then
reduced
under
conditions
Scheme 1
Synthesis of 2-M (M = Cr, Fe) from 1-M, KC8 and 2.2.2-
cryptand.
a Chemistry Division, Los Alamos National Laboratory, Los Alamos, NM 87545,
USA. E-mail: kiplinger@lanl.gov, bstein@lanl.gov, cgoodwin@lanl.gov
b National High Magnetic Field Laboratory, Florida State University,
Tallahassee, FL 32310, USA
c Department of Chemistry \& Biochemistry, Florida State University, Tallahassee,
FL 32306, USA
d Department of Chemistry, School of Natural Sciences, The University of
Manchester, Oxford Road, Manchester, M13 9PL, UK
† Electronic supplementary information (ESI) available. CCDC 2003385–2003390,
2017380 and 2017381. For ESI and crystallographic data in CIF or other electronic
format see DOI: 10.1039/d0cc06518h
Received 11th November 2020,
Accepted 30th November 2020
DOI: 10.1039/d0cc06518h
rsc.li/chemcomm
ChemComm
COMMUNICATION
Published on 18 December 2020. Downloaded by Freie Universitaet Berlin on 11/28/2025 2:25:57 PM. 
View Article Online
View Journal  | View Issue

\newpage

% Page 2

596 | Chem. Commun., 2021, 57, 595--598
This journal is © The Royal Society of Chemistry 2021
established for the synthesis of low-valent metallocenes.16,25
Addition of solid KC8 to cooled (35 1C) solutions of 1-M and
2.2.2-crypt in THF caused rapid color changes (from red to red-
brown, 2-Cr; or red to green-black, 2-Fe). Intensely colored crystals
of 2-Cr and 2-Fe were obtained upon workup (Fig. 1 and ESI†).
These specific Cr and Fe complexes (1-Cr and 1-Fe) were
chosen as they have established redox chemistry and thus
provided an ideal testbed for the comparison of spin-states in
reduced metallocenes.22,23 We were unable to synthesize 1-Mn
and 1-Co in order to furnish a redox series from Cr2+ (d4) to
Co1+ (d8). Reactions between Mn2+ halides and KTMP gave
intractable mixtures. For Co2+ we could isolate traces of impure
[Co(Z5-TMP)(m:Z1:Z1-TMP)]2 (7) (ESI†).
Salient structural parameters for 1-M are shown along with
those of 2-M in Table 1. Complexes 2-M (M = Cr, Fe) are
isomorphous, their ATR-FTIR spectra are superimposable,
and elemental analysis results were in excellent agreement with
prediction. The structures contain discrete [K(2.2.2-crypt)]
cations, and [M(TMP)2] anions. In 2-M the TMP-P atoms are
each situated opposite a b-carbon of the second ring, similar to
neutral [Fe(PC4H2-2,5-Me2)2].26 Upon reduction there is a small
decrease in M–P distances for Cr (D = 0.022(1) Å) but an increase
for Fe (D = +0.103(1) Å); and changes in the TMPcentM distances
follow the same trend (Cr, D = 0.057 Å; Fe, D = +0.053 Å). The
TMPcentCr change between d4 1-Cr (1.795(1) Å) and d5 2-Cr
(avg. 1.738(1) Å) is similar to the Cp*centMn change between d4
and d5 manganocenes.27 The TMPcentFe change between d6
1-Fe (1.660(1) Å) and d7 2-Fe (avg. 1.713(1) Å) is larger than the
increase from 3-Fe (avg. 1.715(2) Å) to 4-Fe (avg. 1.750(3) Å), likely
as the distances in 3-Fe are due to the bulky ligands rather than
electronic effects.16
The 1H, 13C\{1H\}, and 31P\{1H\} NMR spectra of 2-M were
largely uninformative (ESI†) due to paramagnetic broadening
and diamagnetic impurities which dominate the spectra. How-
ever, in the 1H NMR spectrum of 2-Fe we observed three broad
peaks (n1/2 o 17 Hz) in the range 3.37–3.30 ppm corresponding
to 2.2.2-crypt; a doublet at 2.19 ppm (3JHP = 10.29 Hz) and a
singlet at 1.94 ppm correspond to TMP ligand resonances.
Complex 2-Cr could be heated in THF solution (up to 50 1C)
and recrystallized with 480\% recovery of material, whereas
THF solutions of 2-Fe show decomposition in the 1H NMR
spectrum over the course of an hour at room temperature. Data
for 1-M have been reported previously.
We have employed SQUID magnetometry to measure the
susceptibility of the open-shell complexes. 1-Cr and 2-Cr exhibit
wT (298 K) values of B1.24 and B0.38 emu K mol1 (Fig. S34,
ESI†) which suggest S = 1 as previously reported,21b and
S =
1
2 respectively. The observed wT (298 K) value of 2-Fe
(B0.26 emu K mol1) is less than the spin-only value for an
S = 1
2 system (0.375 emu K mol1) (Fig. 2 and Fig. S35, ESI†),
which we attribute to some thermal decomposition of para-
magnetic 2-Fe to diamagnetic 1-Fe, as well as weighing errors
and
diamagnetic
corrections
affecting
this
weakly
para-
magnetic system. As magnetic measurements were inconclusive
for 2-Fe, we turned to high-field electron paramagnetic resonance
(EPR) spectroscopy to determine the ground state spin of 2-Cr
and 2-Fe. These measurements reveal spectra that are typical of
complexes where S = 1
2. The spectrum of 2-Cr is nearly axial with
g1 = 2.023(5), g2 = 1.993(5), and g3 = 1.985(5) (Fig. S32, ESI†); and
2-Fe exhibits a rhombic signal (Fig. 2b) with g1 = 2.033(5),
g2 = 1.999(5), and g3 = 1.943(5), which is similar to other low-
spin Fe1+ sandwich complexes.17
Fig. 1
Molecular structure of 2-Fe with selective atom labelling (C = grey,
P = purple, K = violet, Fe = orange, N = blue, O = red). Displacement
ellipsoids set at 50\% probability level and H-atoms omitted for clarity.
Table 1
Selected structural parameters for 1-M and 2-M
1-Cr
2-Cr
1-Fe
2-Fe
P–M/Å
2.3812(4)
2.3596(8)
2.2932(4)
2.3842(8)
2.3595(7)
2.4078(7)
TMPcent  M/Å
1.795(1)
1.739(1)
1.660(1)
1.714(1)
1.737(1)
1.711(1)
PC2plane  C4plane/1
4.45(12)
3.57(2)
1.27(13)
6.58(2)
4.64(2)
8.73(2)
TMP-P2 twist angle/1
180
139.82(5)
180
146.35(5)
Fig. 2
(a) Temperature dependence of wT for 2-Fe; (b) experimental
(black) and simulated (red) EPR spectra of 2-Fe recorded at 118 GHz and
5 K. The spectrum is simulated with g1 = 2.033, g2 = 1.999, and g3 = 1.943.
Baseline features in the 2-Fe spectrum result from polycrystallinity.
A minor signal arising from an unidentified impurity at lower field is omitted
(see Fig. S28, ESI†); (c) experimental (black) and simulated (red)
57Fe
Mo¨ssbauer spectra of 1-Fe and 2-Fe. Simulations were generated for
1-Fe using d = 0.48(3), DEQ = 2.02(2) mm s1, and for 2-Fe with
d = 0.65(4), DEQ = 1.28(4) mm s1.
Communication
ChemComm
Published on 18 December 2020. Downloaded by Freie Universitaet Berlin on 11/28/2025 2:25:57 PM. 
View Article Online

\newpage

% Page 3

This journal is ©The Royal Society of Chemistry 2021
Chem. Commun., 2021, 57, 595--598 | 597
57Fe Mo¨ssbauer spectroscopy was used to further investigate
the differences between 1-Fe and 2-Fe which display a single
quadrupole doublet in each spectrum (at 120 K, Fig. 2c). The
parameters for 1-Fe, 2-Fe, 3-Fe and 4-Fe are summarized in
Table 2. The more positive isomer shift (d) of 2-Fe suggests
that there is decreased s-orbital electron density at the Fe
nucleus compared to that in 1-Fe. The parameters of 1-Fe are
similar to those reported for derivatized ferrocenes.28 d for 2-Fe
(0.65(4) mm s1) is slightly more positive than many characterized
Fe(I) complexes (B0.2 to B0.6 mm s1),29 but in good agreement
with the low-spin 19 VE [Fe(CpR)(arene)] complexes,17,30 and in
contrast to high-spin 4-Fe (1.25 mm s1).16 We also found excellent
agreement between the DFT calculated d and DEQ parameters and
the experimental data (Table 2) for both 1-Fe and 2-Fe. Data for
3-Fe and 4-Fe is shown for completeness.
To rationalize the structural features of 2-M with their
ground-state
electronic
structures,
we
have
performed
CASSCF/NEVPT2 calculations (see ESI† for details). The ab
initio ligand field theory (AILFT) derived molecular orbital
(MO) diagrams for 1-M, 2-M, 3-Fe, and 4-Fe, are shown in
Fig. 3.16 The most striking observation is the significant differ-
ence in ligand field strength between 2-Fe and 4-Fe, which
results in a drastic difference between the energetic separation
of the dz2 and dxz/dyz orbitals and is responsible for the change
in spin-state of 2-Fe compared to 4-Fe. This can be rationalized
by comparing the Lcent  Fe distances (L = TMP, avg. 1.713(1) Å;
L = Cpttt, avg. 1.750(3) Å); the closer TMP ligand provides a
stronger ligand field than Cpttt, similar to the situation in high-
spin/low-spin manganocenes.31
By considering the MO diagram of a 1,10-diphosphame-
tallocene,32 which is similar to that of FcH, we can rationalize some
of the bond changes.33 In the MO scheme of D5d symmetric FcH the
LUMO is comprised of a doubly degenerate set of antibonding
orbitals with ligand e00
1 (for a D5h Cp ring) contributions that mix with
the dxz/dyz (e1g) orbitals. The LUMO of 1,10-diphosphametallocenes
is also comprised of metal dxz/dyz orbitals, along with a single ligand
HOMO (pP here, Fig. S36, ESI†) that approximates the Cp e00
1 MOs by
having a single nodal plane perpendicular to the ring plane, but has
significant contribution from the P-atom. The SOMO of 2-Fe is
principally dxz in character, and is anti-bonding with respect to pP
which explains the increased Fe–P and TMPcentFe distances.
Conversely, as the SOMO for 2-Cr is comprised of the non-
bonding dz2 orbital as it is in 1-Cr, the structural changes on
reduction are minimal. The slight shortening of TMPcentCr from
1-Cr to 2-Cr is related to the change in spin state from S = 1 to S = 1
2
which lowers the energy of the dx2y2, dxy orbitals, increasing
covalency as in [Cr(Cp)2].34
We suggest that the removal of degeneracy in the ligand
HOMO orbitals for TMP vs. the equivalent orbitals in Cp lifts
the near-degeneracy between the dxz/dyz orbitals (as in 2-Fe
compared to 4-Fe); and furthermore, the large change in ligand
steric profile has drastically changed the separation between dz2
and the dxz/dyz orbitals. These combined effects have stabilized
the low-spin state in 2-Fe and should prove instructive towards
chemical control of the spin state in reduced metallocenes and
heterometallocenes for the 3d series. For example, by signifi-
cantly changing the gap between the dz2 and dxz/dyz orbitals in
4-Fe, a low-spin ferrocenate might be realized.
To summarize, we have demonstrated that octamethyl-1,10-
diphosphametallocenes for Cr and Fe can be reduced to afford
crystalline anionic complexes. Comparison of 2-Fe to authentic
formal Fe(I) complexes suggests that it also contains a formal
Fe(I). We targeted complexes that might be temperature stable;
while 2-Fe is somewhat thermally unstable, 2-Cr appears stable
indefinitely in the solid-state at room temperature. Both were
more amenable to characterization than 4-Fe which decom-
poses in the solid-state even at 35 1C.16 We have found that
both anionic species studied here exhibit low spin (S = 1
2)
ground states, in contrast to high-spin (S = 3
2) 4-Fe, and that
these results are well explained by changes in the ligand
frontier orbitals and M–TMP distances. This suggests that
synthetic control of the ground spin state of 3d metallocenates
is possible and that with judicious planning, it should be
possible to target desired high- or low-spin examples of this
family.
SMG acknowledges support from a Director’s Postdoctoral
Fellowship (20180759PRD4), and CAPG was sponsored by a
Table 2
Mo¨ssbauer parameters for 1-Fe to 4-Fe
1-Fe
2-Fe
3-Feb
4-Feb
d (mm s1)
Exp.:
0.48(3)
0.65(4)
0.66(2)
1.25(2)
Theorya:
0.50
0.66
0.64
1.09
DEQ (mm s1)
Exp.:
2.02(2)
1.28(4)
2.60(2)
1.23(2)
Theorya:
1.91
1.46
2.52
0.65
a Calculated with BP86 and Fe (CP(PPP)), C (def2-tzvp), H (def2-svp)
basis set combination. This gave the best agreement with previous
experimental results for complexes 3-Fe and 4-Fe. b From ref. 16.
Fig. 3
AILFT orbital energies for the frontier orbitals with principally
d-character for 1-M, 2-M, 3-Fe and 4-Fe.16 The indicated electronic
configuration corresponds to the dominant configuration of the CASSCF +
NEVPT2 calculated ground state, and all six complexes feature the same
orbital ordering (from low- to high-energy: dx2y2, dxy, dz2, dxz, dyz). The
molecular z-axis was chosen to align approximately with the vector formed
by the metal and ligand ring-centroids. The SOMO orbital(s) of the reduced
anionic complexes has been indicated for 2-M and 4-Fe.
ChemComm
Communication
Published on 18 December 2020. Downloaded by Freie Universitaet Berlin on 11/28/2025 2:25:57 PM. 
View Article Online

\newpage

% Page 4

598 | Chem. Commun., 2021, 57, 595--598
This journal is © The Royal Society of Chemistry 2021
J. Robert. Oppenheimer Distinguished Postdoctoral Fellowship
(20180703PRD1), both through the Los Alamos National
Laboratory – Laboratory Directed Research and Development
program (LANL-LDRD). CAPG also thanks the UK Engineering
and Physical Sciences Research Council (EPSRC) for a Doctoral
Prize Fellowship (UoM). An EPSRC grant (EP/K039547/1) pro-
vided a single-crystal X-ray diffractometer (UoM). JLK also
thanks the LANL-LDRD program. We also acknowledge the
University of California and LANL Postdoctoral Entrepreneurial
Fellowship and the Feynman Center for Innovation for support
of RJB. BWS acknowledges funding provided by the Director,
Office of Science, Office of Basic Energy Sciences, Division
of Chemical Sciences, Geosciences, and Biosciences, Heavy
Element Chemistry Program of the U.S. Department of Energy
(DOE). We thank Prof. Michael Shatruk for supervising collec-
tion of magnetometry data; and Dr David P. Mills, Dr Fabrizio
Ortu, Dr Andrew J. Gaunt and Dr Michael T. Janicke for helpful
discussions. We also thank Prof. James M. Boncella for the
provision of synthetic work space. A portion of this work was
supported by the NSF award CHE-1955754, and performed at
the National High Magnetic Field Laboratory under National
Science Foundation Cooperative Agreement No. DMR-1644779
with the State of Florida. LANL, an affirmative action/equal
opportunity employer, is managed by Triad National Security,
LLC,
for
the
NNSA
of
the
U.S.
Department
of
Energy
(89233218CNA000001). Raw research data files supporting this
publication are available from Figshare at DOI: 10.6084/
m9.figshare.13348460. The ESI† contains ORCA input files
used in this work.
CAPG lead the project and performed all synthesis, UV-vis-
nIR and NMR characterization; and also X-ray diffraction
measurements with supervision by BLS. SMG performed EPR
and Mo¨ssbauer spectroscopies and interpretation and all the-
oretical work, with supervision by BWS. OU performed SQUID
magnetometry measurements and interpretation. RJB sealed
samples for analysis and was supervised by JLK. The manu-
script was prepared by CAPG with input from all authors.
Conflicts of interest
There are no conflicts to declare.
Notes and references
1 N. J. Long, Metallocenes: An Introduction to Sandwich Complexes,
Wiley, 1998.
2 A. Togni and R. L. Halterman, Metallocenes: Synthesis Reactivity
Applications, 1998.
3 (a) T. J. Kealy and P. L. Pauson, Nature, 1951, 168, 1039; (b) S. A. Miller,
J. A. Tebboth and J. F. Tremaine, J. Chem. Soc., 1952, 632; (c) J. D. Dunitz,
L. E. Orgel and A. Rich, Acta Crystallogr., 1956, 9, 373.
4 G. Wilkinson, M. Rosenblum, M. C. Whiting and R. B. Woodward,
J. Am. Chem. Soc., 1952, 74, 2125.
5 Sˇ. Toma and R. Sˇebesta, Synthesis, 2015, 1683.
6 R. R. Gagne, C. A. Koval and G. C. Lisensky, Inorg. Chem., 1980,
19, 2854.
7 (a) N. Kuhn, E. M. Horn, R. Boese and N. Augart, Angew. Chem., Int.
Ed. Engl., 1988, 27, 1368; (b) O. J. Scherer and T. Bru¨ck, Angew.
Chem., Int. Ed. Engl., 1987, 26, 59; (c) M. Saito, N. Matsunaga,
J. Hamada, S. Furukawa, T. Tada and R. H. Herber, Chem. Lett.,
2019, 48, 163; (d) F. Mathey, Coord. Chem. Rev., 1994, 137, 1.
8 M. Walczak, K. Walczak, R. Mink, M. D. Rausch and G. Stucky, J. Am.
Chem. Soc., 1978, 100, 6382.
9 W. Erb and F. Mongin, Synthesis, 2019, 146.
10 B. Deschamps, J. Fischer, F. Mathey and A. Mitschler, Inorg. Chem.,
1981, 20, 3252.
11 B. Deschamps, J. Fischer, F. Mathey, A. Mitschler and L. Ricard,
Organometallics, 1982, 1, 312.
12 M. Malischewski, M. Adelhardt, J. Sutter, K. Meyer and K. Seppelt,
Science, 2016, 353, 678.
13 (a) A. J. Bard, E. Garcia, S. Kukharenko and V. V. Strelets, Inorg.
Chem., 1993, 32, 3528; (b) Y. Mugnier, C. Moise, J. Tirouflet and
E. Laviron, J. Organomet. Chem., 1980, 186, C49; (c) W. E. Geiger,
J. Am. Chem. Soc., 1974, 96, 2632; (d) J. D. L. Holloway, W. L. Bowden
and W. E. Geiger, J. Am. Chem. Soc., 1977, 99, 7089; (e) N. Ito, T. Saji
and S. Aoyagui, J. Organomet. Chem., 1983, 247, 301.
14 (a) M. Malischewski and K. Seppelt, Dalton Trans., 2019, 48, 17078;
(b) J. L. Robbins, N. M. Edelstein, S. R. Cooper and J. C. Smart, J. Am.
Chem. Soc., 1979, 101, 3853; (c) J. C. Smart and J. L. Robbins, J. Am.
Chem. Soc., 1978, 100, 3936.
15 F. Hung-Low and C. A. Bradley, Inorg. Chem., 2013, 52, 2446.
16 C. A. P. Goodwin, M. Giansiracusa, S. M. Greer, H. M. Nicholas,
P. Evans, M. Vonci, S. Hill, N. F. Chilton and D. P. Mills, Nat. Chem.,
2020, DOI: 10.1038/s41557-020-00595-w.
17 (a) D. Astruc, Tetrahedron, 1983, 39, 4027; (b) M. Rajasekharan,
S. Giezynski, J. Ammeter, N. Oswald, J. Hamon, D. Astruc and
P. Michaud, J. Am. Chem. Soc., 1982, 104, 2400.
18 D. Baudry and M. Ephritikhine, J. Organomet. Chem., 1980, 195,
213.
19 B. M. Gardner, J. McMaster, W. Lewis and S. T. Liddle, Chem.
Commun., 2009, 2851.
20 (a) R. Loschen, C. Loschen, W. Frank and C. Ganter, Eur. J. Inorg.
Chem., 2007, 553; (b) F. Mathey, J. Organomet. Chem., 2002, 646, 15;
(c) M. Ogasawara, T. Nagano, K. Yoshida and T. Hayashi, Organo-
metallics,
2002,
21,
3062;
(d)
M.
Ogasawara,
M.
Shintani,
S. Watanabe, T. Sakamoto, K. Nakajima and T. Takahashi, Organo-
metallics, 2011, 30, 1487.
21 (a) Y. Cabon, D. Carmichael and L. Ricard, Chem. Commun., 2011,
47, 11486; (b) P. Lemoine, M. Gross, P. Braunstein, F. Mathey,
B. Deschamps and J. H. Nelson, Organometallics, 1984, 3, 1303;
(c) N. G. Connelly and W. E. Geiger, Chem. Rev., 1996, 96, 877.
22 R. Feher, F. H. Ko¨hler, F. Nief, L. Ricard and S. Rossmayer, Organo-
metallics, 1997, 16, 4606.
23 (a) A. J. Ashe III, S. Al-Ahmad, S. Pilotek, D. B. Puranik,
C.
Elschenbroich
and
A.
Behrendt,
Organometallics,
1995,
14, 2689; (b) F. Nief, F. Mathey, L. Ricard and F. Robert, Organome-
tallics, 1988, 7, 921.
24 1-Cr was reported vs. the [Co(Cp)2]+/0 couple, and corrected to values
vs. FcH+/0.
25 W. J. Evans, Organometallics, 2016, 35, 3088.
26 G. De Lauzon, B. Deschamps, J. Fischer, F. Mathey and A. Mitschler,
J. Am. Chem. Soc., 1980, 102, 994.
27 (a) N. Augart, R. Boese and G. Schmid, Z. Anorg. Allg. Chem., 1991,
595, 27; (b) W. E. Broderick, J. A. Thompson, E. P. Day and
B. M. Hoffman, Science, 1990, 249, 401; (c) D. P. Freyberg,
J. L. Robbins, K. N. Raymond and J. C. Smart, J. Am. Chem. Soc.,
1979, 101, 892.
28 (a) R. M. G. Roberts, J. Silver and B. Yamin, J. Organomet. Chem.,
1984, 270, 221; (b) R. A. Stukan, S. P. Gubin, A. N. Nesmeyanov,
V. I. Gol’danskii and E. F. Makarov, Theor. Exp. Chem., 1966, 2, 581;
(c) G. K. Wertheim and R. H. Herber, J. Chem. Phys., 1963, 38,
2106.
29 J. P. Mariot, P. Michaud, S. Lauer, D. Astruc, A. X. Trautwein and
F. Varret, J. Phys., 1983, 44, 1377.
30 P. Gu¨tlich, E. Bill and A. X. Trautwein, Mo¨ssbauer Spectroscopy and
Transition Metal Chemistry, 2011.
31 M. D. Walter, C. D. Sofield, C. H. Booth and R. A. Andersen,
Organometallics, 2009, 28, 2005.
32 N. M. Kostic and R. F. Fenske, Organometallics, 1983, 2, 1008.
33 A. Haaland, Acc. Chem. Res., 1979, 12, 415.
34 M. Swart, Inorg. Chim. Acta, 2007, 360, 179.
Communication
ChemComm
Published on 18 December 2020. Downloaded by Freie Universitaet Berlin on 11/28/2025 2:25:57 PM. 
View Article Online

\newpage

\end{document}